% SPDX-License-Identifier: 0BSD
\section{Reusable Leading-Colour Execution Layouts}
\label{sec:lc-flow-layouts}

Leading-colour coverage and its numerical layout are separate concepts.  Let
\(\mathcal H\) be the retained helicities and \(\mathcal F\) the retained
physical colour flows.  Complete leading-colour coverage represents
\begin{equation}
  \mathcal C_{\mathrm{LC}}=\mathcal H\times\mathcal F
  \label{eq:lc-complete-coverage}
\end{equation}
and can expose every resolved amplitude \(\mathcal A_{hf}\).  For selected
subsets \(S_H\subseteq\mathcal H\) and \(S_F\subseteq\mathcal F\), the
observable is
\begin{equation}
  \mathcal W(S_H,S_F)
  =\mathcal N\sum_{h\in S_H}\sum_{f\in S_F}w_{hf}
   \left|\sum_{r\in G(h,f)}\mathcal A^{(r)}_{hf}\right|^2 .
  \label{eq:lc-runtime-selection}
\end{equation}
Here \(G(h,f)\) is a coherent contribution group, \(w_{hf}\) contains exact
symmetry and leading-colour weights, and \(\mathcal N\) is the process
normalization.  The layout changes how this sum is scheduled, not its
components or normalization.

\subsection{Two complete-coverage layouts}

The report measures two layouts because selected-flow and all-flow workloads
offer different opportunities for sharing:

\begin{center}
\small
\begin{tabular}{@{}L{1.32in}L{2.40in}L{2.40in}@{}}
\toprule
\textbf{property} & \textbf{topology replay} &
  \textbf{all-flow union} \\
\midrule
Measured workload & one selected flow, summed helicities &
  all flows, one selected helicity \\
Retained physics & every flow and helicity &
  every flow and helicity \\
Main reuse & equivalent flow topologies related by exact permutations and
  factors & currents and closures shared directly across all flows \\
Availability & compiled, eager, recurrence, and on-the-fly to
  \(n_{\mathrm{final}}=4\) &
  compiled, eager, and recurrence \\
\bottomrule
\end{tabular}
\end{center}

On-the-fly execution does not materialize a separate all-flow-union artifact.
Within its LC scope, one compact topology-replay artifact serves both measured
workloads: one selected flow with helicities summed, and all physical flows
with one helicity selected. Thus the all-flow observable remains the same, but
its on-the-fly execution is reconstructed lazily from the shared compact
artifact rather than from \(D_{\mathrm{union}}\).

In topology replay, exact relations partition physical flows into equivalence
classes. If \(\rho(f)\) represents the class containing \(f\), the calculation
records a permutation \(\pi_f\), an exact factor \(s_f\), and the corresponding
weight:
\begin{equation}
  \mathcal A_{h f}(p)
  =s_f\,\mathcal A_{\pi_f(h),\rho(f)}\!\left(\pi_f(p)\right).
  \label{eq:lc-topology-replay}
\end{equation}
This avoids constructing an independent current graph for every physical
ordering and is suited to selected-flow, helicity-summed evaluation.

The all-flow union instead merges the dependency graphs of all resolved
components:
\begin{equation}
  D_{\mathrm{union}}
  =\left(\bigcup_{h\in\mathcal H,\,f\in\mathcal F}D_{hf}\right)
   \big/\!\sim ,
  \label{eq:lc-all-flow-union}
\end{equation}
where \(\sim\) identifies currents, kernel values, and closures that are
algebraically equal up to retained exact factors.  For a selected helicity,
one graph traversal then reuses intermediate currents across all flows.

Compiled mode specializes the chosen layout to process-wide stage evaluators.
Eager mode applies prepared-model kernels to the same dependency structure,
recurrence mode constructs a compact current schedule directly, and on-the-fly
mode reconstructs query-local traces from its compact process seed. These
execution choices do not change \(\mathcal H\), \(\mathcal F\), or
Eq.~\eqref{eq:lc-runtime-selection}.

\subsection{Runtime selection}

Flow and helicity selections refer to physical components.  A selection may be
applied to a complete batch or specified independently for each point; omitting
an axis sums all retained entries on that axis.  Rows with equal selections can
be grouped into contiguous numerical batches.  This grouping changes neither
the generated coverage nor the resolved result.

\subsection{Generation and reference boundaries}

The two layouts build different graphs and numerical schedules, so their
process-generation times are distinct observables:
\begin{equation}
  T_{\mathrm{gen}}^{\mathrm{replay}}
  \quad\hbox{and}\quad
  T_{\mathrm{gen}}^{\mathrm{union}}.
  \label{eq:lc-two-generation-times}
\end{equation}
Cross-mode comparisons use the same layout on both sides.  Preparation of
reusable model-level kernels occurs before the process-generation timer and is
not included in these table values.

The on-the-fly exception has one process-generation quantity \(G\), because
the selected-flow and all-flow workloads share one compact artifact. Its first
full-batch evaluation \(W\) constructs the initially requested selector-family
state. Artifact loading and the conventional benchmark warmup calls are timed
separately and are excluded from \(W\). Calling \texttt{Runtime.clear()} drops
the lazy on-the-fly state without replacing the process artifact.

The original-\AC{} reference uses two different routes. For the selected-flow
workload, its generation value measures creation of the process library. For
the all-flow workload, the reported setup value measures preparation of direct
all-flow evaluation. The corresponding runtime measurements implement the
same physical workloads as \PAC{}, but the all-flow setup quantities do not
define a like-for-like process-generation speedup. The table therefore reports
both quantities absolutely and marks their generation comparison as not
comparable.

Before timing an on-the-fly cell, the complete resolved components for both
workloads are compared with separately generated recurrence and compiled
authorities. The same componentwise comparisons are repeated after
\texttt{Runtime.clear()} to check the cold reconstruction path.

\subsection{NLC and full colour}

The alternative layouts apply only at leading colour.  NLC and full-colour
evaluation construct an amplitude basis and contract it with the appropriate
colour metric:
\begin{equation}
  \mathcal W_{\mathrm{contracted}}
  =\mathcal N\sum_h\sum_{a,b}
   \mathcal A_{ha}^{*}\,C_{ab}\,\mathcal A_{hb}.
  \label{eq:contracted-colour-layout}
\end{equation}
These treatments do not expose the physical leading-colour flow axis, so there
is no topology-replay or all-flow-union distinction. On-the-fly NLC and
full-colour entries are explicit future placeholders rather than runnable
measurements.
